% Options for packages loaded elsewhere
\PassOptionsToPackage{unicode}{hyperref}
\PassOptionsToPackage{hyphens}{url}
%
\documentclass[
  11pt,
  letterpaper,
]{article}

\usepackage{amsmath,amssymb}
\usepackage{setspace}
\usepackage{iftex}
\ifPDFTeX
  \usepackage[T1]{fontenc}
  \usepackage[utf8]{inputenc}
  \usepackage{textcomp} % provide euro and other symbols
\else % if luatex or xetex
  \usepackage{unicode-math}
  \defaultfontfeatures{Scale=MatchLowercase}
  \defaultfontfeatures[\rmfamily]{Ligatures=TeX,Scale=1}
\fi
\usepackage{lmodern}
\ifPDFTeX\else  
    % xetex/luatex font selection
\fi
% Use upquote if available, for straight quotes in verbatim environments
\IfFileExists{upquote.sty}{\usepackage{upquote}}{}
\IfFileExists{microtype.sty}{% use microtype if available
  \usepackage[]{microtype}
  \UseMicrotypeSet[protrusion]{basicmath} % disable protrusion for tt fonts
}{}
\makeatletter
\@ifundefined{KOMAClassName}{% if non-KOMA class
  \IfFileExists{parskip.sty}{%
    \usepackage{parskip}
  }{% else
    \setlength{\parindent}{0pt}
    \setlength{\parskip}{6pt plus 2pt minus 1pt}}
}{% if KOMA class
  \KOMAoptions{parskip=half}}
\makeatother
\usepackage{xcolor}
\usepackage[top=1in,bottom=1in,left=1in,right=1in]{geometry}
\setlength{\emergencystretch}{3em} % prevent overfull lines
\setcounter{secnumdepth}{-\maxdimen} % remove section numbering
% Make \paragraph and \subparagraph free-standing
\makeatletter
\ifx\paragraph\undefined\else
  \let\oldparagraph\paragraph
  \renewcommand{\paragraph}{
    \@ifstar
      \xxxParagraphStar
      \xxxParagraphNoStar
  }
  \newcommand{\xxxParagraphStar}[1]{\oldparagraph*{#1}\mbox{}}
  \newcommand{\xxxParagraphNoStar}[1]{\oldparagraph{#1}\mbox{}}
\fi
\ifx\subparagraph\undefined\else
  \let\oldsubparagraph\subparagraph
  \renewcommand{\subparagraph}{
    \@ifstar
      \xxxSubParagraphStar
      \xxxSubParagraphNoStar
  }
  \newcommand{\xxxSubParagraphStar}[1]{\oldsubparagraph*{#1}\mbox{}}
  \newcommand{\xxxSubParagraphNoStar}[1]{\oldsubparagraph{#1}\mbox{}}
\fi
\makeatother


\providecommand{\tightlist}{%
  \setlength{\itemsep}{0pt}\setlength{\parskip}{0pt}}\usepackage{longtable,booktabs,array}
\usepackage{calc} % for calculating minipage widths
% Correct order of tables after \paragraph or \subparagraph
\usepackage{etoolbox}
\makeatletter
\patchcmd\longtable{\par}{\if@noskipsec\mbox{}\fi\par}{}{}
\makeatother
% Allow footnotes in longtable head/foot
\IfFileExists{footnotehyper.sty}{\usepackage{footnotehyper}}{\usepackage{footnote}}
\makesavenoteenv{longtable}
\usepackage{graphicx}
\makeatletter
\def\maxwidth{\ifdim\Gin@nat@width>\linewidth\linewidth\else\Gin@nat@width\fi}
\def\maxheight{\ifdim\Gin@nat@height>\textheight\textheight\else\Gin@nat@height\fi}
\makeatother
% Scale images if necessary, so that they will not overflow the page
% margins by default, and it is still possible to overwrite the defaults
% using explicit options in \includegraphics[width, height, ...]{}
\setkeys{Gin}{width=\maxwidth,height=\maxheight,keepaspectratio}
% Set default figure placement to htbp
\makeatletter
\def\fps@figure{htbp}
\makeatother

\usepackage{setspace}
\usepackage{float}
\usepackage{microtype}
\usepackage{fancyhdr}
\usepackage{lastpage}
\usepackage{enumitem}
\usepackage{booktabs}
\usepackage{longtable}
\usepackage{array}
\usepackage{amsmath,amssymb,amsthm}
\usepackage{graphicx}
\usepackage{xurl}
\usepackage[hidelinks]{hyperref}

\setlength{\parindent}{0pt}
\setlength{\parskip}{4pt}
\setlength{\textfloatsep}{8pt plus 2pt minus 2pt}
\setlength{\floatsep}{6pt plus 2pt minus 2pt}
\setlength{\intextsep}{8pt plus 2pt minus 2pt}

\pagestyle{fancy}
\fancyhf{}
\fancyfoot[C]{\thepage}
\renewcommand{\headrulewidth}{0pt}
\renewcommand{\footrulewidth}{0pt}
\makeatletter
\@ifpackageloaded{caption}{}{\usepackage{caption}}
\AtBeginDocument{%
\ifdefined\contentsname
  \renewcommand*\contentsname{Table of contents}
\else
  \newcommand\contentsname{Table of contents}
\fi
\ifdefined\listfigurename
  \renewcommand*\listfigurename{List of Figures}
\else
  \newcommand\listfigurename{List of Figures}
\fi
\ifdefined\listtablename
  \renewcommand*\listtablename{List of Tables}
\else
  \newcommand\listtablename{List of Tables}
\fi
\ifdefined\figurename
  \renewcommand*\figurename{Figure}
\else
  \newcommand\figurename{Figure}
\fi
\ifdefined\tablename
  \renewcommand*\tablename{Table}
\else
  \newcommand\tablename{Table}
\fi
}
\@ifpackageloaded{float}{}{\usepackage{float}}
\floatstyle{ruled}
\@ifundefined{c@chapter}{\newfloat{codelisting}{h}{lop}}{\newfloat{codelisting}{h}{lop}[chapter]}
\floatname{codelisting}{Listing}
\newcommand*\listoflistings{\listof{codelisting}{List of Listings}}
\makeatother
\makeatletter
\makeatother
\makeatletter
\@ifpackageloaded{caption}{}{\usepackage{caption}}
\@ifpackageloaded{subcaption}{}{\usepackage{subcaption}}
\makeatother

\ifLuaTeX
  \usepackage{selnolig}  % disable illegal ligatures
\fi
\usepackage[]{natbib}
\bibliographystyle{plainnat}
\usepackage{bookmark}

\IfFileExists{xurl.sty}{\usepackage{xurl}}{} % add URL line breaks if available
\urlstyle{same} % disable monospaced font for URLs
\hypersetup{
  pdftitle={Stat: Penalized Zero-Inflated Beta Regression},
  pdfauthor={Swarnali Banerjee and Gregory J. Matthews},
  hidelinks,
  pdfcreator={LaTeX via pandoc}}


\title{Stat: Penalized Zero-Inflated Beta Regression}
\author{Swarnali Banerjee and Gregory J. Matthews}
\date{2026-09-17}

\begin{document}
\maketitle


\setstretch{1}
\section{Project Summary}\label{project-summary}

\subsection{Overview}\label{overview}

{[}State the scientific problem in one or two sentences.{]}

{[}Explain why the problem matters and what is currently missing.{]}

{[}State the central research objective and the overall strategy.{]}

The central hypothesis is that \textbf{{[}CENTRAL HYPOTHESIS{]}}.

The project will pursue three complementary aims:

\begin{enumerate}
\def\labelenumi{\arabic{enumi}.}
\tightlist
\item
  \textbf{Aim 1:} {[}FOUNDATIONAL THEORY OR METHOD{]}
\item
  \textbf{Aim 2:} {[}STATISTICAL METHODOLOGY{]}
\item
  \textbf{Aim 3:} {[}VALIDATION, COMPUTATION, OR SCIENTIFIC
  APPLICATION{]}
\end{enumerate}

The expected outcome is {[}MAIN CONTRIBUTION{]}.

\subsection{Intellectual Merit}\label{intellectual-merit}

{[}Explain the fundamental statistical/mathematical knowledge that will
be advanced.{]}

{[}Describe what is genuinely new relative to existing theory and
methods.{]}

{[}State the expected theoretical, methodological, or computational
contributions.{]}

\subsection{Broader Impacts}\label{broader-impacts}

{[}Describe concrete impacts through training, software, data,
education, interdisciplinary research, dissemination, broadening
participation, or other appropriate mechanisms.{]}

{[}State how impacts will be assessed.{]}

\section{Project Description}\label{project-description}

\subsection{1. Motivation and
Background}\label{motivation-and-background}

Proportion data arise in many scientific settings. Examples 1, 2, 3 etc.
Traditionally this type of data has modeled using normal regressions
techniques with transformations of the response variable such as the the
arcsin square root transformation. However, this method is known to be
flawed. Some have even referred to it as ``asinine''. Other alternatives
include using logistic/ binomial regression, however, these require the
statistical practitioner to know the denominator (i.e.~the \(n\) of
these proportions). In many settings all we have is the proportion. In
these cases, a better alternative is beta regression (CITATION). Beta
regression is well studied (IS IT?) with established packages available
in R. \citet{Kosmidis2026}

https://pubmed.ncbi.nlm.nih.gov/21560670/

{[}Introduce the broad scientific or statistical problem.{]}

{[}Explain the current state of knowledge.{]}

{[}Identify the specific limitation or gap.{]}

For example, existing approaches typically rely on {[}ASSUMPTION{]},
which is problematic when {[}SETTING{]}. Existing methods such as
{[}METHODS{]} provide useful solutions, but do not address {[}SPECIFIC
GAP{]} \citep{example2024}.

The proposed research addresses this gap by developing {[}NEW
FRAMEWORK{]}.

\subsection{2. Central Research
Question}\label{central-research-question}

The central research question is:

\begin{quote}
\textbf{{[}CENTRAL RESEARCH QUESTION{]}}
\end{quote}

The project is guided by the following hypothesis:

\begin{quote}
\textbf{{[}CENTRAL HYPOTHESIS{]}}
\end{quote}

\subsection{3. Overall Strategy}\label{overall-strategy}

The proposed work combines {[}THEORY{]}, {[}STATISTICAL METHODOLOGY{]},
and {[}COMPUTATION/APPLICATION{]}.

The aims are deliberately complementary. Aim 1 establishes
{[}FOUNDATION{]}; Aim 2 develops {[}METHOD{]}; and Aim 3 evaluates and
demonstrates {[}CONTRIBUTION{]}.

The project is designed so that failure of any single technical approach
will not prevent progress toward the overall objective.

\subsection{4. Research Plan}\label{research-plan}

\subsubsection{4.1 Specific Aims}\label{specific-aims}

\begin{longtable}[]{@{}lll@{}}
\toprule\noalign{}
Aim & Scientific objective & Primary deliverable \\
\midrule\noalign{}
\endhead
\bottomrule\noalign{}
\endlastfoot
1 & {[}OBJECTIVE{]} & {[}THEOREM/METHOD{]} \\
2 & {[}OBJECTIVE{]} & {[}METHOD/SOFTWARE{]} \\
3 & {[}OBJECTIVE{]} & {[}VALIDATION/APPLICATION{]} \\
\end{longtable}

\subsubsection{4.2 Conceptual Framework}\label{conceptual-framework}

The proposed framework is:

The framework links {[}COMPONENT 1{]} to {[}COMPONENT 2{]} through
{[}MECHANISM{]}.

\subsubsection{4.3 Evaluation Strategy}\label{evaluation-strategy}

Each aim will be evaluated using a combination of theoretical analysis,
simulation, computational benchmarking, and/or scientific validation.

Success will be assessed using prespecified criteria:

\begin{itemize}
\tightlist
\item
  {[}CRITERION 1{]}
\item
  {[}CRITERION 2{]}
\item
  {[}CRITERION 3{]}
\item
  {[}CRITERION 4{]}
\end{itemize}

\subsection{5. Aim 1: {[}FOUNDATIONAL THEORY OR
METHOD{]}}\label{aim-1-foundational-theory-or-method}

\subsubsection{Motivation}\label{motivation}

{[}Explain the specific gap addressed by Aim 1.{]}

\subsubsection{Research Question}\label{research-question}

{[}Specific question.{]}

\subsubsection{Proposed Approach}\label{proposed-approach}

We will develop {[}METHOD{]}.

Let \(Y\) denote {[}OUTCOME{]}, \(X\) denote {[}COVARIATES{]}, and
\(\theta\) denote {[}PARAMETERS{]}. We consider the general model

\[
[MODEL]
\]

The principal theoretical objectives are:

\begin{enumerate}
\def\labelenumi{\arabic{enumi}.}
\tightlist
\item
  {[}OBJECTIVE{]}
\item
  {[}OBJECTIVE{]}
\item
  {[}OBJECTIVE{]}
\end{enumerate}

\subsubsection{Theoretical Development}\label{theoretical-development}

{[}Describe identifiability.{]}

{[}Describe consistency.{]}

{[}Describe asymptotic behavior.{]}

{[}Describe efficiency, robustness, or other relevant properties.{]}

\subsubsection{Simulation Study}\label{simulation-study}

The simulation study will vary:

\begin{longtable}[]{@{}ll@{}}
\toprule\noalign{}
Factor & Levels \\
\midrule\noalign{}
\endhead
\bottomrule\noalign{}
\endlastfoot
Sample size & {[}LEVELS{]} \\
Signal strength & {[}LEVELS{]} \\
Missingness/dependence structure & {[}LEVELS{]} \\
Model misspecification & {[}LEVELS{]} \\
\end{longtable}

Primary performance measures will include {[}BIAS{]}, {[}RMSE{]},
{[}COVERAGE{]}, {[}POWER{]}, and {[}COMPUTATIONAL COST{]}.

\subsubsection{Expected Outcomes}\label{expected-outcomes}

Aim 1 will establish {[}EXPECTED CONTRIBUTION{]}.

\subsubsection{Potential Problems and Alternative
Strategies}\label{potential-problems-and-alternative-strategies}

{[}Technical risk.{]}

{[}Alternative strategy.{]}

{[}Criteria for adopting the alternative.{]}

\subsection{6. Aim 2: {[}STATISTICAL
METHODOLOGY{]}}\label{aim-2-statistical-methodology}

\subsubsection{Motivation}\label{motivation-1}

{[}Explain why Aim 1 is insufficient by itself and what methodological
problem remains.{]}

\subsubsection{Proposed Method}\label{proposed-method}

{[}Detailed description of estimator/model/algorithm.{]}

The proposed estimator is

\[
\widehat{\theta} = [FORMULA].
\]

{[}Explain estimation.{]}

{[}Explain uncertainty quantification.{]}

{[}Explain computational implementation.{]}

\subsubsection{Benchmarking}\label{benchmarking}

The proposed method will be compared descriptively with established
alternatives including {[}METHOD 1{]}, {[}METHOD 2{]}, and {[}METHOD
3{]}.

Evaluation will consider statistical performance, robustness,
interpretability, and computational scalability.

\subsubsection{Software}\label{software}

The methods will be implemented in an open-source R package, tentatively
named \texttt{{[}PACKAGE\ NAME{]}}.

The package will include:

\begin{itemize}
\tightlist
\item
  reproducible implementations of all proposed methods;
\item
  simulation utilities;
\item
  diagnostics;
\item
  documentation and examples; and
\item
  vignettes reproducing the principal analyses.
\end{itemize}

\subsubsection{Expected Outcomes}\label{expected-outcomes-1}

{[}Expected contribution.{]}

\subsubsection{Potential Problems and Alternative
Strategies}\label{potential-problems-and-alternative-strategies-1}

{[}Potential problem and alternative.{]}

\subsection{7. Aim 3: {[}VALIDATION OR SCIENTIFIC
APPLICATION{]}}\label{aim-3-validation-or-scientific-application}

\subsubsection{Scientific Context}\label{scientific-context}

{[}Describe the scientific problem for which the proposed methodology is
relevant.{]}

\subsubsection{Data}\label{data}

{[}Describe data source, sample size, variables, and study design.{]}

\subsubsection{Analysis}\label{analysis}

{[}Specify primary analysis.{]}

{[}Specify secondary analyses.{]}

{[}Specify sensitivity analyses.{]}

\subsubsection{Validation}\label{validation}

{[}Describe external validation, simulation-based calibration, held-out
prediction, replication, or other validation strategy.{]}

\subsubsection{Expected Outcomes}\label{expected-outcomes-2}

{[}Expected scientific contribution.{]}

\subsubsection{Integration Across Aims}\label{integration-across-aims}

{[}Explain how Aim 3 tests whether the theory and methodology developed
in Aims 1 and 2 produce useful scientific knowledge.{]}

\subsection{8. Preliminary Work}\label{preliminary-work}

Our preliminary work establishes feasibility of the proposed research.

\subsubsection{Preliminary Result 1}\label{preliminary-result-1}

{[}Describe preliminary analysis.{]}

\subsubsection{Preliminary Result 2}\label{preliminary-result-2}

{[}Describe second result.{]}

\begin{longtable}[]{@{}lrrr@{}}
\toprule\noalign{}
Method & Bias & RMSE & Coverage \\
\midrule\noalign{}
\endhead
\bottomrule\noalign{}
\endlastfoot
{[}Existing method{]} & {[} {]} & {[} {]} & {[} {]} \\
\hyperref[proposed-method]{Proposed method} & {[} {]} & {[} {]} & {[}
{]} \\
\end{longtable}

These preliminary results motivate the proposed development of
{[}METHOD{]}.

\subsection{9. Intellectual Merit}\label{intellectual-merit-1}

The intellectual merit of the project rests on {[}CENTRAL SCIENTIFIC
CONTRIBUTION{]}.

\subsubsection{Fundamental Contribution}\label{fundamental-contribution}

{[}Explain the fundamental statistical or mathematical question.{]}

\subsubsection{Methodological
Innovation}\label{methodological-innovation}

The proposed work advances current methodology by:

\begin{enumerate}
\def\labelenumi{\arabic{enumi}.}
\tightlist
\item
  {[}INNOVATION{]}
\item
  {[}INNOVATION{]}
\item
  {[}INNOVATION{]}
\end{enumerate}

\subsubsection{Generalizability}\label{generalizability}

{[}Explain why the results will apply beyond the motivating
application.{]}

The expected contributions include {[}THEORY{]}, {[}METHODOLOGY{]},
{[}COMPUTATION{]}, and \hyperref[software]{SOFTWARE}.

\subsection{10. Broader Impacts}\label{broader-impacts-1}

\subsection{Broader Impacts}\label{broader-impacts-2}

Chat GPT write this: The proposed research will develop statistical
methodology for analyzing bounded continuous outcomes that exhibit
substantial excess zeros. Such outcomes arise in a wide range of
scientific and social-scientific settings, including proportions, rates,
utilization measures, adherence measures, ecological indices, and other
outcomes for which zero represents a meaningful boundary value rather
than a missing observation. The proposed development of zero-inflated
penalized beta regression will provide researchers with a principled
framework for modeling these data while accommodating both the boundary
mass at zero and the continuous variation among positive observations.

A primary broader impact of the project will be to improve the
accessibility and reliability of statistical methods for investigators
who encounter these data in applied research. Existing approaches often
require researchers to choose between models that inadequately represent
the zero-generating process and models that fail to account for the
bounded nature of the nonzero observations. By developing estimation
procedures, inferential methods, diagnostics, and computational tools
specifically for zero-inflated beta regression with penalization, this
project will provide a practical alternative for researchers working
with high-dimensional or otherwise challenging datasets. The resulting
methods will be designed for implementation in commonly used statistical
software, lowering the barrier between methodological development and
applied use.

The project will also contribute to statistical education and training.
Modern applied statisticians increasingly encounter outcomes that are
simultaneously bounded, zero-inflated, and associated with many
potential predictors. These characteristics are often treated separately
in introductory statistical modeling, despite their frequent occurrence
together in contemporary datasets. The proposed work will provide
examples and computational materials that can be incorporated into
graduate courses in regression, generalized linear models,
high-dimensional statistics, and statistical computing. Reproducible
examples will illustrate how model assumptions, regularization, zero
inflation, and bounded outcomes interact, providing students with a
framework for thinking critically about model selection rather than
treating individual modeling techniques as isolated procedures.

An additional impact will be the development and dissemination of
open-source computational resources. Software implementing the proposed
methodology, together with documentation, examples, simulation code, and
reproducible analyses, will be made publicly available. These resources
will facilitate independent evaluation of the methods, encourage
adoption by applied researchers, and provide a foundation for subsequent
methodological extensions. Whenever possible, computational
implementations will use established statistical software ecosystems so
that the methods can be incorporated into existing analytical workflows.

The project will further support the development of researchers through
mentoring and training in statistical methodology, computation, and
reproducible research. Students involved in the project will gain
experience translating theoretical statistical developments into
computational algorithms and practical analytical tools. Training will
emphasize not only mathematical and computational skills but also
communication of statistical methodology to researchers outside
statistics. This combination of methodological depth and applied
communication is increasingly important as statistical models are used
across scientific disciplines.

Finally, the proposed research has the potential to broaden the role of
statistical methodology in scientific investigations where conventional
regression models do not adequately reflect the structure of the
observed data. Rather than developing a method for a single application,
the project addresses a general class of statistical problems. By making
the resulting methodology theoretically justified, computationally
accessible, and reproducible, the project will contribute to the broader
statistical infrastructure available to researchers studying bounded and
zero-inflated outcomes.

\subsubsection{10.1 Training}\label{training}

{[}Describe graduate and undergraduate training.{]}

Students will receive training in:

\begin{itemize}
\tightlist
\item
  statistical theory;
\item
  reproducible research;
\item
  R programming;
\item
  scientific communication; and
\item
  collaborative research.
\end{itemize}

\subsubsection{10.2 Open Science and
Software}\label{open-science-and-software}

{[}Describe planned software and code release.{]}

All appropriate software will be released under an open-source license,
with documentation and reproducible examples.

\subsubsection{10.3 Education and
Dissemination}\label{education-and-dissemination}

{[}Courses, workshops, public-facing materials, conferences, tutorials,
etc.{]}

\subsubsection{10.4 Broadening
Participation}\label{broadening-participation}

{[}Concrete activities, partnerships, mentoring, or recruitment
strategies.{]}

\subsubsection{10.5 Assessment}\label{assessment}

\begin{longtable}[]{@{}
  >{\raggedright\arraybackslash}p{(\columnwidth - 4\tabcolsep) * \real{0.3333}}
  >{\raggedright\arraybackslash}p{(\columnwidth - 4\tabcolsep) * \real{0.3333}}
  >{\raggedright\arraybackslash}p{(\columnwidth - 4\tabcolsep) * \real{0.3333}}@{}}
\toprule\noalign{}
\begin{minipage}[b]{\linewidth}\raggedright
Activity
\end{minipage} & \begin{minipage}[b]{\linewidth}\raggedright
Metric
\end{minipage} & \begin{minipage}[b]{\linewidth}\raggedright
Target
\end{minipage} \\
\midrule\noalign{}
\endhead
\bottomrule\noalign{}
\endlastfoot
Student training & Students completing research experience &
{[}TARGET{]} \\
Software & Downloads/citations/users & {[}TARGET{]} \\
Dissemination & Presentations/workshops & {[}TARGET{]} \\
Publications & Peer-reviewed products & {[}TARGET{]} \\
\end{longtable}

\subsection{11. Timeline and Project
Management}\label{timeline-and-project-management}

The project will be conducted over four years.

\begin{longtable}[]{@{}lcccc@{}}
\toprule\noalign{}
Activity & Year 1 & Year 2 & Year 3 & Year 4 \\
\midrule\noalign{}
\endhead
\bottomrule\noalign{}
\endlastfoot
Aim 1 & X & X & & \\
Aim 2 & & X & X & \\
Aim 3 & & & X & X \\
Software & X & X & X & X \\
Student training & X & X & X & X \\
Dissemination & & X & X & X \\
\end{longtable}

\subsubsection{Year 1}\label{year-1}

{[}Milestones.{]}

\subsubsection{Year 2}\label{year-2}

{[}Milestones.{]}

\subsubsection{Year 3}\label{year-3}

{[}Milestones.{]}

\subsubsection{Year 4}\label{year-4}

{[}Milestones.{]}

\subsubsection{Risk Management}\label{risk-management}

{[}Describe dependencies and contingency plans.{]}

\subsection{12. Results from Prior NSF
Support}\label{results-from-prior-nsf-support}

\subsubsection{Award {[}AWARD NUMBER{]}}\label{award-award-number}

\textbf{Title:} {[}TITLE{]}\\
\textbf{Amount:} \${[}AMOUNT{]}\\
\textbf{Period:} {[}START--END{]}

\paragraph{Intellectual Merit}\label{intellectual-merit-2}

{[}Results.{]}

\paragraph{Broader Impacts}\label{broader-impacts-3}

{[}Results.{]}

\paragraph{Publications and Products}\label{publications-and-products}

{[}Publications, software, datasets, presentations, or other
products.{]}

\paragraph{Relationship to Proposed
Work}\label{relationship-to-proposed-work}

{[}Explain relationship to the current project.{]}

\subsubsection{Award {[}SECOND AWARD, IF
APPLICABLE{]}}\label{award-second-award-if-applicable}

{[}Repeat structure.{]}

If no qualifying prior NSF support exists, replace this section with the
appropriate NSF-required statement.

\section{References Cited}\label{references-cited}

\renewcommand{\bibsection}{}
\bibliography{references.bib}





\end{document}
